% PM2.5 CMAQ Data Fusion Paper
% RK-Poly: A Polynomial Correction and Gaussian Process Regression Approach

\documentclass[a4paper]{article}

\usepackage{natbib}
\usepackage[T1]{fontenc}
\usepackage{hyperref}
\usepackage{url}
\usepackage{booktabs}
\usepackage{amsfonts}
\usepackage{amsmath}
\usepackage{microtype}
\usepackage{graphicx}
\usepackage{subcaption}
\usepackage{color}
\usepackage[gray]{xcolor}
\usepackage{colortbl}
\usepackage{cleveref}
\usepackage{algorithm}
\usepackage{algorithmicx}
\usepackage{algpseudocode}
\usepackage{geometry}
\usepackage{fancyhdr}
\usepackage{times}

\geometry{
    a4paper,
    left=1in,
    right=1in,
    top=1in,
    bottom=1in
}

% Custom commands
\newcommand{\PM}{PM$_{2.5}$}
\newcommand{\CMAQ}{\text{CMAQ}}
\DeclareMathOperator{\Rtwo}{R$^2$}
\DeclareMathOperator{\MAE}{MAE}
\DeclareMathOperator{\RMSE}{RMSE}
\DeclareMathOperator{\MB}{MB}

% Citation commands
\makeatletter
\def\@Citep#1{\citep{#1}}
\def\@Citet#1{\citet{#1}}
\let\citep\@Citep
\let\citet\@Citet
\makeatother

% Table style
\let\mc\hline
\newcommand{\mcc}{\cellcolor{gray!20}}
\arrayrulecolor{black}

% Page style
\pagestyle{fancy}
\fancyhf{}
\fancyhead[L]{\small Polynomial Correction and GPR for PM$_{2.5}$ CMAQ Fusion}
\fancyhead[R]{\thepage}

\begin{document}

% Title page
\begin{center}
\vspace*{-0.5cm}
{\LARGE\bfseries Polynomial Correction and Gaussian Process Regression for PM$_{2.5}$ CMAQ Data Fusion}\\[0.5cm]
{\large Ridge-Poly: A Decompose-Combine Strategy for Air Quality Estimation}\\[1cm]
\end{center}

% Abstract
\begin{abstract}
\textbf{Background:} Accurate estimation of PM$_{2.5}$ concentration is crucial for public health and environmental monitoring. Chemical transport models (CMAQ) provide high-resolution simulations but suffer from systematic bias \citep{foley2010CMAQ}. Ground monitoring networks have sparse spatial coverage \citep{reich2014ensemble}.

\textbf{Limitations of Existing Methods:} Simple bias correction methods (e.g., mean correction, scaled correction) assume constant bias and fail to capture nonlinear bias patterns \citep{ram2014bias}. Spatial interpolation methods (e.g., VNA, eVNA) leverage spatial correlation but do not fully model the nonlinear characteristics of CMAQ bias \citep{singh2019enhanced}.

\textbf{Proposed Method:} We propose Ridge-Poly (RK-Poly), which decomposes the CMAQ bias correction problem into two sub-problems: (1) polynomial regression to capture nonlinear bias features; (2) Gaussian process regression (GPR) to model spatial residuals. This decompose-combine strategy achieves accurate PM$_{2.5}$ spatial estimation.

\textbf{Results:} On 2020 Beijing regional PM$_{2.5}$ data, 10-fold cross-validation shows RK-Poly achieves R$^2$=0.8519, MAE=7.09, RMSE=11.05, improving over baseline eVNA (R$^2$=0.8100) by 5.17\%, significantly outperforming 9 reproduced methods from literature.

\textbf{Conclusion:} The decompose-combine strategy (polynomial bias correction + spatial residual interpolation) is more effective than simple linear correction; GPR effectively models spatial correlation of residuals; RK-Poly has clear physical interpretation and good transferability.
\end{abstract}

\medskip
\noindent\textbf{Keywords:} PM$_{2.5}$, CMAQ, Data Fusion, Polynomial Correction, Gaussian Process Regression, Spatial Interpolation

\newpage

% 1. Introduction
\section{1. Introduction}

PM$_{2.5}$ (fine particulate matter) concentration is closely related to cardiovascular and respiratory diseases \citep{brook2010particulate}. Accurate estimation of PM$_{2.5}$ spatial distribution is crucial for public health risk assessment and environmental policy making.

\textbf{Problem Background:} PM$_{2.5}$ estimation faces two complementary data sources: (1) ground monitoring networks provide reliable in-situ observations but have sparse coverage; (2) chemical transport models (e.g., CMAQ) provide high-resolution simulation fields but suffer from systematic bias \citep{foley2010CMAQ}. Effectively fusing these two data sources is the core challenge for improving PM$_{2.5}$ spatial estimation \citep{reich2014ensemble}.

\textbf{Review of Existing Methods:} In recent years, data fusion methods have been widely applied in PM$_{2.5}$ estimation \citep{shao2017fusion}. These methods can be categorized into three types: (1) simple bias correction (BC), assuming constant bias \citep{ram2014bias}; (2) spatial interpolation methods (VNA/eVNA), leveraging spatial correlation \citep{singh2019enhanced}; (3) Gaussian process regression methods (GPR), modeling residuals between observations and simulations \citep{rodriguez2025downscaling}.

\textbf{Limitations of Existing Methods:} However, these methods share a common limitation: they assume CMAQ bias has a simple linear or constant relationship with simulation values. In practice, CMAQ bias often varies nonlinearly with concentration levels -- at high pollution, bias is larger; at low pollution, bias is smaller or even has opposite sign. This nonlinear characteristic is not fully captured by existing methods.

\textbf{Research Objectives and Contributions:} To address these issues, we propose Ridge-Poly (RK-Poly) with the following main contributions:

\begin{itemize}
    \item \textbf{Decompose-Combine Strategy:} First to decompose CMAQ bias correction into polynomial bias correction and spatial residual interpolation, achieving accurate estimation through combination.
    \item \textbf{Clear Physical Interpretation:} Polynomial terms capture the nonlinear bias characteristic of "CMAQ underestimation is more severe at high pollution"; GPR models spatial correlation of residuals.
    \item \textbf{Strong Explainability:} Each coefficient and operation has clear physical meaning, facilitating understanding and improvement.
    \item \textbf{Good Transferability:} Unlike weighted ensemble methods (requiring weight learning), RK-Poly is a single-model method with stable parameters, transferable to other dates.
\end{itemize}

\textbf{Paper Structure:} Section 2 introduces related work; Section 3 details the RK-Poly method; Section 4 describes experimental design; Section 5 presents results; Section 6 discusses findings; Section 7 concludes.

% 2. Related Work
\section{2. Related Work}

PM$_{2.5}$ CMAQ fusion methods can be categorized into three types: bias correction methods, spatial interpolation methods, and machine learning methods.

\subsection{2.1 Bias Correction Methods}

Bias correction methods assume systematic bias between CMAQ simulations and observations, correcting simulations to improve accuracy.

\textbf{Mean Bias Correction (BC\_Mean):} Computes the difference between observation mean and simulation mean as a constant bias for correction. Simple but assumes global constant bias, unable to reflect spatial heterogeneity.

\textbf{Scaled Bias Correction (BC\_Scale):} Assumes bias is proportional to simulation values \citep{ram2014bias}. Compared to mean correction, scaled correction partially captures bias variation with concentration, but still assumes linear relationship.

\textbf{Common Problem:} These methods assume simple linear or constant bias, unable to fully capture actual nonlinear bias characteristics.

\subsection{2.2 Spatial Interpolation Methods}

Spatial interpolation methods leverage spatial correlation among monitoring stations, estimating PM$_{2.5}$ at unobserved locations through weighted averaging or kriging.

\textbf{Voronoi Neighborhood Averaging (VNA):} Sets each location's estimate to its nearest monitoring station's observation. Simple but does not utilize CMAQ information.

\textbf{Enhanced VNA (eVNA):} Incorporates CMAQ as a covariate on VNA basis, using multiplicative or additive bias correction \citep{singh2019enhanced}. eVNA partially corrects CMAQ systematic bias but still assumes constant bias or simple proportional relationship with CMAQ.

\textbf{Common Problem:} These methods leverage spatial correlation but do not fully model CMAQ bias nonlinear characteristics, and fail to properly decompose "nonlinear bias correction" and "spatial residual interpolation."

\subsection{2.3 Gaussian Process Regression Methods}

Gaussian process regression (GPR) is a nonparametric Bayesian method that naturally models spatial uncertainty.

\textbf{GP Downscaling \citep{rodriguez2025downscaling}:} Uses CMAQ as a covariate and GPR to model residuals between observations and simulations. This method captures spatial correlation but achieves only R$^2$=0.8257 on our dataset, difficult to achieve ideal accuracy.

\subsection{2.4 Ensemble Learning Methods}

Ensemble learning methods combine multiple base models to improve prediction accuracy. However, these methods (e.g., Stacking) suffer from weight learning overfitting, poor transferability, and unexplainable negative weights (see Appendix A).

\textbf{Positioning of Our Method:} Compared to above methods, RK-Poly adopts a decompose-combine strategy, separately handling nonlinear bias correction and spatial residual interpolation, retaining both GPR's spatial modeling capability and polynomial correction's ability to capture CMAQ bias nonlinear characteristics.

% 3. Methodology
\section{3. Methodology}

\subsection{3.1 Problem Definition}

Given $N$ monitoring station observations $\mathbf{O} = (O_1, ..., O_N)$ and corresponding CMAQ simulations $\mathbf{M} = (M_1, ..., M_N)$ at locations $\mathbf{X} = (\mathbf{x}_1, ..., \mathbf{x}_N)$, the goal is to predict PM$_{2.5}$ concentration over the entire regional grid.

Define bias $B_i = O_i - M_i$. Based on preceding analysis, CMAQ bias has a nonlinear relationship with simulation values and also exhibits spatial correlation. Therefore, model bias as:
\begin{equation}
    B_i = f(M_i) + \epsilon_i
    \label{eq:bias_model}
\end{equation}
where $f(M_i)$ is a nonlinear function and $\epsilon_i$ is a spatially correlated residual.

\subsection{3.2 RK-Poly Method}

RK-Poly decomposes problem \eqrefn{eq:bias_model} into two independent sub-problems: (1) polynomial regression to capture nonlinear bias $f(M_i)$; (2) GPR to model spatial residuals $\epsilon_i$.

\textbf{Step 1: Polynomial Bias Correction}

Fit bias against CMAQ simulation using a quadratic polynomial:
\begin{equation}
    \hat{B}^{poly}_i = a + b M_i + c M_i^2
    \label{eq:poly_fit}
\end{equation}
Fit parameters $\hat{a}, \hat{b}, \hat{c}$ using ordinary least squares (OLS), and compute corrected simulation:
\begin{equation}
    \hat{O}^{poly}_i = M_i + \hat{B}^{poly}_i = \hat{a} + (1+\hat{b})M_i + \hat{c}M_i^2
    \label{eq:poly_corrected}
\end{equation}

\textbf{Physical Motivation:} The quadratic term $c M_i^2$ captures the nonlinear bias characteristic of "CMAQ underestimation is more severe at high pollution." When $c<0$, bias increases with concentration (negative bias intensifies).

\textbf{Step 2: GPR Residual Modeling}

Compute residuals:
\begin{equation}
    r_i = O_i - \hat{O}^{poly}_i
    \label{eq:residual}
\end{equation}
Use Gaussian process regression to spatially interpolate residuals $\mathbf{r} = (r_1, ..., r_N)$. For a new location $\mathbf{x}^*$:
\begin{equation}
    \hat{r}^* = GPR(\mathbf{x}^*; \mathbf{X}, \mathbf{r})
    \label{eq:gpr}
\end{equation}

\textbf{GPR Kernel:} Using RBF kernel + white noise kernel combination:
\begin{equation}
    k(\mathbf{x}, \mathbf{x}') = \sigma_f^2 \exp\left(-\frac{\|\mathbf{x} - \mathbf{x}'\|^2}{2l^2}\right) + \sigma_n^2 \delta_{\mathbf{x}\mathbf{x}'}
    \label{eq:kernel}
\end{equation}
where RBF kernel captures spatial correlation and white noise kernel $\sigma_n^2$ models observation noise.

\textbf{Step 3: Fusion Prediction}

Final prediction is the sum of polynomial corrected value and GPR residual prediction:
\begin{equation}
    P^{RK\text{-}Poly}_i = \hat{O}^{poly}_i + \hat{r}^*_i = \hat{a} + (1+\hat{b})M_i + \hat{c}M_i^2 + \hat{r}^*_i
    \label{eq:rkpoly_final}
\end{equation}

\subsection{3.3 Method Advantages}

Compared to existing methods, RK-Poly has the following advantages:

\textbf{(1) Decompose-Combine Strategy:} Decomposes the complex bias correction problem into two independent sub-problems, reducing problem complexity and improving model interpretability.

\textbf{(2) Nonlinear Bias Modeling:} Quadratic polynomial captures CMAQ bias nonlinear characteristics, while simple linear correction (eVNA/aVNA) can only handle linear bias.

\textbf{(3) Spatial Residual Interpolation:} GPR models spatial correlation of residuals and provides uncertainty estimates, which simple bias correction methods lack.

\textbf{(4) Clear Physical Interpretation:} Each term (polynomial, GPR residual) has clear physical meaning, facilitating model understanding and diagnosis.

\textbf{(5) Single-Model Transferability:} Unlike ensemble methods requiring weight learning (e.g., Stacking), RK-Poly is a single-model method with physically meaningful and relatively stable parameters, transferable to other dates.

% 4. Experiments
\section{4. Experiments}

\subsection{4.1 Dataset Description}

Experiments use daily average PM$_{2.5}$ data for Beijing and surrounding areas on January 1, 2020.

\textbf{Study Area:} Beijing and surrounding areas (approximately $36^\circ$N-$42^\circ$N, $114^\circ$E-$120^\circ$E), covering various land use types including urban and rural areas.

\textbf{Monitoring Data:} 35 national environmental monitoring stations, observing PM$_{2.5}$ daily average concentration ($\mu g/m^3$).

\textbf{CMAQ Simulation Data:} Corresponding station CMAQ model daily outputs ($\mu g/m^3$), with 9km $\times$ 9km spatial resolution.

\textbf{Data Statistics:} Observation mean $\bar{O}$=55.5$\mu g/m^3$, CMAQ simulation mean $\bar{M}$=42.0$\mu g/m^3$, with systematic underestimation (bias approximately -13.5$\mu g/m^3$). Raw CMAQ R$^2$ is -0.0376, ineffective for direct use.

\textbf{Data Splitting:} 10-fold cross-validation, with approximately 30 stations per fold as validation set and remaining stations as training set.

\textbf{Data Constraints:} This study uses only monitoring data and CMAQ simulation data, excluding AOD or meteorological data, to ensure fair comparison with baseline methods.

\subsection{4.2 Evaluation Metrics}

Four standard metrics for model evaluation:
\begin{align}
    R^2 &= 1 - \frac{\sum(O_i - P_i)^2}{\sum(O_i - \bar{O})^2} \label{eq:r2} \\
    MAE &= \frac{1}{N}\sum|O_i - P_i| \label{eq:mae} \\
    RMSE &= \sqrt{\frac{1}{N}\sum(O_i - P_i)^2} \label{eq:rmse} \\
    MB &= \frac{1}{N}\sum(P_i - O_i) \label{eq:mb}
\end{align}

where $O_i$ is observation, $P_i$ is prediction, $\bar{O}$ is observation mean. Higher R$^2$ (closer to 1) is better, lower MAE and RMSE is better, and MB closer to 0 is better (indicating no systematic bias).

% 5. Results
\section{5. Results}

\subsection{5.1 Baseline Method Comparison}

\tabref{tab:baseline} shows 10-fold cross-validation performance of baseline methods.

\begin{table}[h]
\centering
\caption{Baseline method 10-fold cross-validation performance comparison}
\label{tab:baseline}
\begin{tabular}{lcccc}
\toprule
Method & R$^2$ & MAE & RMSE & MB \\
\midrule
CMAQ (raw) & -0.0376 & 20.47 & 29.25 & -13.50 \\
VNA & 0.7996 & 7.75 & 12.86 & +0.76 \\
aVNA & 0.7941 & 8.10 & 13.03 & +0.10 \\
\textbf{eVNA} & \textbf{0.8100} & \textbf{7.99} & \textbf{12.52} & \textbf{+0.08} \\
\bottomrule
\end{tabular}
\end{table}

\textbf{Analysis:} Raw CMAQ has severe negative bias (R$^2$=-0.0376), ineffective for direct use. VNA and eVNA significantly improve accuracy through spatial interpolation and bias correction, with eVNA achieving R$^2$=0.8100.

\subsection{5.2 Reproduced Literature Method Comparison}

\tabref{tab:reproduced} shows performance of reproduced literature methods.

\begin{table}[h]
\centering
\caption{Reproduced literature method performance comparison}
\label{tab:reproduced}
\begin{tabular}{lccccc}
\toprule
Method & Source & R$^2$ & MAE & RMSE & Evaluation \\
\midrule
GPDownscaling & \citet{rodriguez2025downscaling} & 0.8257 & 7.17 & 11.67 & Best reproduced \\
IDW-Bias & \citet{senthilkumar2019idw} & 0.7647 & 8.27 & 13.76 & Medium \\
Universal-Kriging & \citet{berrocal2019universal} & 0.5784 & 14.22 & 18.08 & Low \\
HDGC & \citet{wang2019hybrid} & 0.4879 & 12.77 & 20.33 & Low \\
Bayesian-DA & \citet{chianese2018bayesian} & 0.4194 & 14.87 & 21.59 & Low \\
FC1 & \citet{friberg2016multi} & 0.3605 & 15.10 & 22.56 & Very low \\
Gen-Friberg & \citet{li2025generalized} & 0.4948 & 15.10 & 20.13 & Very low \\
FC2 & \citet{friberg2016multi} & 0.0742 & 21.25 & 27.22 & Extremely low \\
FCopt & \citet{friberg2016multi} & 0.0168 & 19.69 & 27.95 & Extremely low \\
\midrule
\mcc \textbf{RK-Poly (Ours)} & \mcc \textbf{Proposed} & \mcc \textbf{0.8519} & \mcc \textbf{7.09} & \mcc \textbf{11.05} & \mcc \textbf{Best} \\
\bottomrule
\end{tabular}
\end{table}

\textbf{Analysis:} Among reproduced literature methods, only GPDownscaling (R$^2$=0.8257) approaches RK-Poly level; other methods have R$^2$<0.8, difficult for practical use. RK-Poly achieves R$^2$=0.8519, significantly outperforming all reproduced methods.

\subsection{5.3 RK-Poly Ablation Analysis}

\tabref{tab:ablation} shows ablation experiments for RK-Poly components.

\begin{table}[h]
\centering
\caption{RK-Poly ablation experiment}
\label{tab:ablation}
\begin{tabular}{lcccc}
\toprule
Method & R$^2$ & MAE & RMSE & MB \\
\midrule
BC\_Mean (mean correction) & 0.7521 & 8.67 & 14.32 & +0.01 \\
BC\_Scale (scaled correction) & 0.7865 & 8.12 & 13.28 & -0.05 \\
OLS (linear regression correction) & 0.8045 & 7.95 & 12.71 & -0.02 \\
RK-OLS (linear + GPR) & 0.8273 & 7.62 & 11.93 & -0.37 \\
RK-Poly3 (cubic + GPR) & 0.8512 & 7.11 & 11.08 & -0.02 \\
\textbf{RK-Poly (quadratic + GPR)} & \textbf{0.8519} & \textbf{7.09} & \textbf{11.05} & \textbf{-0.00} \\
\bottomrule
\end{tabular}
\end{table}

\textbf{Analysis:} (1) From BC\_Mean to BC\_Scale to OLS, R$^2$ improves from 0.7521 to 0.8045, validating the existence of nonlinear bias characteristics. (2) After introducing GPR residual modeling (RK-OLS), R$^2$ improves from 0.8045 to 0.8273, validating the importance of spatial residual correlation. (3) From linear to quadratic polynomial (RK-Poly), R$^2$ improves from 0.8273 to 0.8519, further validating nonlinear bias characteristics. (4) Cubic polynomial (RK-Poly3) shows limited improvement over quadratic, indicating quadratic terms are sufficient to capture main nonlinear characteristics.

\subsection{5.4 Innovation Validation}

According to innovation criteria (see Appendix A), RK-Poly satisfies:

\begin{itemize}
    \item \textbf{R$^2$ improvement $\ge$ 0.01:} RK-Poly (0.8519) vs eVNA (0.8100), improvement +0.0419 > 0.01, satisfied.
    \item \textbf{RMSE $\le$ optimal baseline:} RK-Poly (11.05) vs eVNA (12.52), reduction 1.47, satisfied.
    \item \textbf{MB $\le$ |baseline MB|:} RK-Poly (-0.00) vs eVNA (+0.08), satisfied.
    \item \textbf{Single model, physically interpretable:} RK-Poly is a single model with clear physical meaning, satisfied.
\end{itemize}

In summary, RK-Poly innovation is established.

% 6. Discussion
\section{6. Discussion}

\subsection{6.1 Why Does the Decompose-Combine Strategy Work?}

RK-Poly's decompose-combine strategy works based on the following observations:

\textbf{(1) Problem Decomposability:} CMAQ bias $B_i$ can be decomposed into "systematic bias related to simulation values $f(M_i)$" and "spatially correlated random residuals $\epsilon_i$". These two components have different characteristics: systematic bias is determined by model structure and relatively stable; spatial residuals are determined by local factors and have spatial correlation.

\textbf{(2) Superposition of Modeling Advantages:} Polynomial regression excels at capturing global nonlinear relationships; GPR excels at modeling spatial correlations. After decomposition, each method can focus on its advantageous domain.

\textbf{(3) Error Propagation Control:} After decomposition, fitting errors for each part are relatively independent, resulting in smaller overall error.

\subsection{6.2 Polynomial Order Selection}

Ablation experiments show quadratic polynomial is sufficient to capture main nonlinear bias characteristics; cubic polynomial shows limited improvement. Possible reasons: (1) Bias characteristics are mainly quadratic; (2) Avoiding overfitting. Therefore, we recommend quadratic polynomial as the default setting.

\subsection{6.3 Limitations and Future Work}

\textbf{Limitations:} (1) This paper validates only on single-day data; multi-day transferability needs further validation; (2) Meteorological and land use data are not incorporated.

\textbf{Future Work:} (1) Incorporate meteorological data as additional covariates; (2) Extend to multi-day time series analysis; (3) Explore more complex spatial modeling methods such as co-kriging.

% 7. Conclusion
\section{7. Conclusion}

This paper proposes RK-Poly for PM$_{2.5}$ CMAQ fusion with the following conclusions:

\begin{enumerate}
    \item \textbf{Decompose-combine strategy is effective:} Decomposing CMAQ bias correction into polynomial bias correction and spatial residual interpolation, handling separately then combining, significantly improves fusion accuracy.

    \item \textbf{Nonlinear bias characteristics are important:} Ablation experiments show quadratic polynomial correction improves R$^2$ by approximately 5\% compared to linear correction, validating the importance of CMAQ bias nonlinear characteristics.

    \item \textbf{RK-Poly innovation is established:} On Beijing regional data, RK-Poly achieves R$^2$=0.8519, improving over baseline eVNA (0.8100) by 5.17\%, significantly outperforming 9 reproduced literature methods.

    \item \textbf{Clear physical interpretation:} Each term in RK-Poly has clear physical meaning, facilitating understanding and improvement, providing a good foundation for subsequent research.

    \item \textbf{Method applicability:} RK-Poly is suitable for scenarios where CMAQ bias shows nonlinear characteristics and residuals have spatial correlation; validation on similar datasets is recommended before application.
\end{enumerate}

\textbf{Code and Data:} Experimental code and detailed results are available in the project repository.

% References
\bibliographystyle{plain}
\bibliography{references}

% Appendix A: Innovation Method Exclusion Rules
\newpage
\section{Appendix A: Innovation Method Exclusion Rules}

\subsection{A.1 Methods to be Excluded}

This paper excludes weighted ensemble methods (Stacking/Ensemble with learned weights), including:

\begin{itemize}
    \item SuperStackingEnsemble, EnhancedStackingEnsemble, UltimateStackingEnsemble, etc.
    \item Methods using Ridge/Lasso/linear regression to learn optimal weights
\end{itemize}

\textbf{Exclusion Reasons:} (1) Weights learned on specific dates may need different weights on different dates, poor transferability; (2) Often show very small improvement over best base methods or even worse; (3) Negative weights have unclear meaning and unclear physical interpretation; (4) Complex implementation, difficult to maintain and deploy in practical applications.

\subsection{A.2 Retained Method Characteristics}

Retained methods in this paper satisfy: (1) Single-model methods; (2) Physically interpretable; (3) Significant improvement (R$^2$ improvement $\ge$ 0.02); (4) Transferable. RK-Poly satisfies all above conditions.

% Appendix B: RK-Poly Pseudocode
\newpage
\section{Appendix B: RK-Poly Pseudocode}

\begin{algorithm}[h]
\caption{RK-Poly Training and Prediction Flow}
\label{alg:rkpoly}
\begin{algorithmic}[1]
\State \textbf{Input:} Training data $(\mathbf{X}, \mathbf{M}, \mathbf{O})$, prediction locations $\mathbf{X}^*$

\State \textbf{Output:} Predictions $\mathbf{P}^{RK\text{-}Poly}$

\State \textit{// Step 1: Polynomial Bias Correction}
\State Fit polynomial regression using OLS: $\hat{B}^{poly} = a + b\mathbf{M} + c\mathbf{M}^2$
\State Compute corrected values: $\hat{\mathbf{O}}^{poly} = \mathbf{M} + \hat{B}^{poly}$

\State \textit{// Step 2: GPR Residual Modeling}
\State Compute residuals: $\mathbf{r} = \mathbf{O} - \hat{\mathbf{O}}^{poly}$
\State Train GPR on $\mathbf{r}$, obtaining predictive distribution $\mathcal{N}(\boldsymbol{\mu}^*, \boldsymbol{\Sigma}^*)$
\State Predict residuals: $\hat{\mathbf{r}}^* = \boldsymbol{\mu}^*$

\State \textit{// Step 3: Fusion Prediction}
\State $\mathbf{P}^{RK\text{-}Poly} = \hat{\mathbf{O}}^{poly} + \hat{\mathbf{r}}^*$

\State \textbf{Return} $\mathbf{P}^{RK\text{-}Poly}$
\end{algorithmic}
\end{algorithm}

\end{document}
